\documentclass[11pt,a4paper]{article}

\usepackage[utf8]{inputenc}
\usepackage[T1]{fontenc}
\usepackage[margin=2.5cm]{geometry}
\usepackage{booktabs}
\usepackage{enumitem}
\usepackage{xcolor}
\usepackage{titlesec}
\usepackage{fancyhdr}
\usepackage[colorlinks=true, linkcolor=red!55!black, urlcolor=red!55!black, citecolor=red!55!black]{hyperref}

\definecolor{wgired}{RGB}{130,20,20}

\titleformat{\section}{\Large\bfseries\color{wgired}}{\thesection}{0.6em}{}
\titleformat{\subsection}{\large\bfseries\color{wgired}}{\thesubsection}{0.6em}{}

\pagestyle{fancy}
\fancyhf{}
\lhead{\small From \texttt{\textbackslash begin} to \texttt{\textbackslash end}}
\rhead{\small White Globe Initiative}
\cfoot{\thepage}
\renewcommand{\headrulewidth}{0.4pt}

\setlist[itemize]{leftmargin=1.4em, itemsep=2pt, topsep=2pt}
\setlist[enumerate]{leftmargin=1.6em, itemsep=2pt, topsep=2pt}
\emergencystretch=2em

\begin{document}

% ---------------- Title Page ----------------
\begin{titlepage}
  \centering
  \vspace*{2.5cm}
  {\Huge\bfseries\color{wgired} From \texttt{\textbackslash begin} to \texttt{\textbackslash end}}\\[10pt]
  {\Large\itshape The Complete LaTeX Guide}\\[6pt]
  {\large Course Outline \& Syllabus}\\[40pt]

  \rule{0.6\textwidth}{0.5pt}\\[20pt]

  \begin{tabular}{rl}
    \textbf{Instructor:}       & Atikur Rahman \\[4pt]
    \textbf{Organized by:}     & White Globe Initiative \\[4pt]
    \textbf{Instructor Email:} & \url{atik@whiteglobeinitiative.org} \\[4pt]
    \textbf{General Contact:}  & \url{contact@whiteglobeinitiative.org} \\[4pt]
    \textbf{Format:}           & 7-Day Intensive Course \\[4pt]
    \textbf{Audience:}         & Complete beginners, economics background \\
  \end{tabular}\\[20pt]

  \rule{0.6\textwidth}{0.5pt}

  \vfill
  {\small A White Globe Initiative in-house training program.}
\end{titlepage}

\tableofcontents
\newpage

% ---------------- Course Overview ----------------
\section{Course Overview}

This course takes a complete beginner from zero LaTeX experience to
independently writing, formatting, citing, and submitting a
publication-ready economics paper in seven days.

\textbf{Audience.} Complete beginners with an economics background,
writing for the purpose of producing publication-quality articles,
working papers, and reports.

\textbf{Format.} Seven daily sessions of 2--3 hours of instruction
plus hands-on practice. Each day builds directly on the last, using a
single running paper draft as the exercise throughout the week.

\textbf{Tooling.} Overleaf (\url{overleaf.com}) -- a free,
browser-based LaTeX editor. No local installation is required, the
free tier is sufficient, and it includes a live PDF preview and a
journal template gallery.

\textbf{Primary references.} Overleaf's ``Learn LaTeX in 30 Minutes,''
the full Overleaf documentation library (\url{overleaf.com/learn}),
and the Wikibooks LaTeX manual (\url{en.wikibooks.org/wiki/LaTeX}) --
used throughout as the topic-by-topic reference set for every day
below.

\textbf{Teaching style.} Every concept in the daily slide decks is
taught as a triad: the \textbf{code} you write, the \textbf{compiled
output} it produces, and a plain-language explanation of \emph{why}
you write it and \emph{what} it does -- a book-style, example-first
approach designed for someone who has never opened a \texttt{.tex}
file before.

\textbf{Outcome.} By Day 7, the trainee independently produces a
formatted, cited, journal-template-ready economics paper and can
self-debug common compile errors without assistance.

% ---------------- Week at a Glance ----------------
\section{Week at a Glance}

\begin{center}
\begin{tabular}{clp{6cm}}
  \toprule
  \textbf{Day} & \textbf{Focus} & \textbf{Key Skill Unlocked} \\
  \midrule
  1 & Setup \& Document Anatomy & Compile a working \texttt{.tex} file from scratch \\
  2 & Math Mode for Economists & Typeset models, equations, and notation \\
  3 & Tables, Figures \& Regression Output & Publication-quality regression tables \\
  4 & Citations \& Bibliography & Author-year citations with BibTeX/natbib \\
  5 & Document Structure \& Long Documents & Organize a full paper across files \\
  6 & Journal Templates \& Formatting & Submission-ready formatting \\
  7 & Troubleshooting \& Independent Mastery & Self-debug and finalize the paper \\
  \bottomrule
\end{tabular}
\end{center}

\newpage
% ---------------- Daily Curriculum ----------------
\section{Daily Curriculum}

\subsection*{Day 1: Setup \& Document Anatomy}
\textbf{Objective.} Get the trainee compiling their first document and
understanding what LaTeX actually is -- a markup language, not a word
processor. Zero prior exposure assumed.

\textbf{Topics.} What LaTeX is and isn't; getting started with
Overleaf; the minimal document (preamble vs.\ body); document class
and packages; the title block (\texttt{\textbackslash title},
\texttt{\textbackslash author}, \texttt{\textbackslash maketitle});
sections and subsections; paragraphs and line breaks; bold, italics,
and comments; reading your first compile errors.

\textbf{Exercise.} Recreate the title page and abstract of a real,
published economics paper, compiling at least five times and
intentionally breaking and fixing something each time.

\textbf{Homework.} Retype the introduction paragraph of the paper you
will rebuild all week into a new Overleaf project.

\vspace{6pt}
\subsection*{Day 2: Math Mode for Economists}
\textbf{Objective.} Typeset models, regression specifications, and
notation with full precision -- no equation editor required.

\textbf{Topics.} Inline vs.\ display math; the \texttt{amsmath}
package; subscripts and superscripts; Greek letters and operators
(sums, products, limits, argmax); fractions and binomials; multi-line
derivations with \texttt{align}; matrices (correlation matrices,
game-theory payoff matrices); labeling and referencing equations.

\textbf{Exercise.} Typeset a Cobb-Douglas production function, a
linear regression equation with a subscripted error term, and a
$2\times2$ payoff matrix.

\textbf{Homework.} Add the full set of equations from the running
paper draft, labeling every numbered equation.

\vspace{6pt}
\subsection*{Day 3: Tables, Figures \& Regression Output}
\textbf{Objective.} Build a repeatable workflow for turning regression
output and charts into clean, journal-style tables and figures.

\textbf{Topics.} The \texttt{tabular} environment; \texttt{booktabs}
rules for the economics-journal look; \texttt{multicolumn}/\texttt{
cmidrule} for grouped regression headers; a complete worked regression
table; the \texttt{figure} environment and \texttt{includegraphics};
float placement; the \url{tablesgenerator.com} workflow for Excel/Stata
output.

\textbf{Exercise.} Reproduce a small regression output table in
booktabs style using \url{tablesgenerator.com} as a first pass, then
hand-edit for captions and labels.

\textbf{Homework.} Insert all tables and figures from the running
paper draft, with proper captions and labels.

\vspace{6pt}
\subsection*{Day 4: Citations \& Bibliography}
\textbf{Objective.} Build a working BibTeX pipeline reusable for every
future paper, practicing on real published research.

\textbf{Topics.} BibTeX basics and the \texttt{.bib} file; \texttt{
natbib} and author-year citation styles common in economics;
\texttt{\textbackslash citet} vs.\ \texttt{\textbackslash citep};
\texttt{biblatex} as an alternative; sourcing \texttt{.bib} entries
from Google Scholar, reference managers, and journal websites;
cross-referencing sections, equations, tables, and figures.

\textbf{Exercise.} Build a five-entry \texttt{.bib} file from real
papers listed at \url{whiteglobeinitiative.org/published-papers}
(or your own field's literature) and write a short mock literature
review citing all five.

\textbf{Homework.} Build the real bibliography for the running paper
draft and cite throughout the introduction/literature review.

\vspace{6pt}
\subsection*{Day 5: Document Structure \& Long-Document Tools}
\textbf{Objective.} Move from a single short file to the structure a
full paper or thesis-length document needs.

\textbf{Topics.} Table of contents; footnotes; the \texttt{hyperref}
package for clickable cross-references; splitting a long paper across
files with \texttt{\textbackslash input}; page numbering, headers, and
footers.

\textbf{Exercise.} Split the running paper draft into separate files
and assemble them into a single \texttt{main.tex} using \texttt{
\textbackslash input}, adding a table of contents and at least one
footnote.

\textbf{Homework.} Load \texttt{hyperref} and confirm every
cross-reference resolves correctly.

\vspace{6pt}
\subsection*{Day 6: Journal Templates \& Professional Formatting}
\textbf{Objective.} Become fluent in adapting an existing journal
template rather than building formatting from scratch.

\textbf{Topics.} Browsing the Overleaf template gallery (AEA-style,
Elsevier's \texttt{elsarticle}, working-paper templates); filling in
title/author/abstract/JEL/keyword metadata; page size and margins;
using colour sparingly; converting an existing Word draft to LaTeX.

\textbf{Exercise.} Migrate the entire running paper draft into one
real economics journal template, end-to-end.

\textbf{Homework.} Proofread the migrated document for formatting
breaks and list everything found.

\vspace{6pt}
\subsection*{Day 7: Troubleshooting, Polish \& Independent Mastery}
\textbf{Objective.} Leave able to diagnose errors and finish a paper
without a tutor in the room.

\textbf{Topics.} Reading compile errors like a pro; search strategy
for TeX StackExchange; the pre-submission proofreading checklist;
exporting the final PDF; where to go next (lshort, the wider Overleaf
library).

\textbf{Exercise.} The instructor injects 3--5 common errors into the
trainee's document; the trainee diagnoses and fixes each one unaided.

\textbf{Final deliverable.} A submission-ready PDF of the paper built
across the week, self-reviewed against the Day 6 template's
requirements.

\newpage
% ---------------- Master Resource List ----------------
\section{Master Resource List}

\begin{center}
\begin{tabular}{p{7cm}p{7.5cm}}
  \toprule
  \textbf{Resource} & \textbf{Best Used For} \\
  \midrule
  Overleaf -- Learn LaTeX in 30 Minutes & Day 1 crash course; first document, structure, math preview \\
  Overleaf Documentation (\url{overleaf.com/learn}) & Topic-by-topic reference for every day of the course \\
  Wikibooks LaTeX (\url{en.wikibooks.org/wiki/LaTeX}) & Free companion textbook; deeper explanations and exercises \\
  tablesgenerator.com & Converting Excel/Stata output into LaTeX table code (Day 3) \\
  whiteglobeinitiative.org/published-papers & Real published-paper examples for citation practice (Day 4) \\
  Overleaf Template Gallery (\url{overleaf.com/latex/templates}) & Economics/AEA/Elsevier journal templates (Day 6) \\
  TeX StackExchange (\url{tex.stackexchange.com}) & Debugging real compile errors (Day 7 and beyond) \\
  ``The Not So Short Introduction to LaTeX2e'' (lshort, free PDF) & Reference manual once basics are fluent \\
  \bottomrule
\end{tabular}
\end{center}

\vspace{20pt}
\begin{center}
\small
White Globe Initiative --- Atikur Rahman\\
\url{atik@whiteglobeinitiative.org} \quad\textbullet\quad \url{contact@whiteglobeinitiative.org}
\end{center}

\end{document}
